% Puissances de 2 et de 10 : le nom de l'unité, la valeur exacte, l'écart.
\documentclass[border=6pt]{standalone}
\usepackage{sd}
\begin{document}
\begin{tikzpicture}[papier, x=21mm]
  \draw[ligne, line width=1pt] (0.6,0) -- (6.4,0);
  \foreach \i/\dix/\deux/\si/\iec/\ecart/\c in {
      1/3/10/Ko/Kio/{2,4}/Pjaune,
      2/6/20/Mo/Mio/{4,9}/Porange,
      3/9/30/Go/Gio/{7,4}/Pbleu,
      4/12/40/To/Tio/{10,0}/Pvert,
      5/15/50/Po/Pio/{12,6}/Plavande,
      6/18/60/Eo/Eio/{15,3}/Prose} {
    \draw[ligne, line width=1pt] (\i,-0.12) -- (\i,0.12);
    \node[bloc=\c, minimum width=17mm, anchor=south] (h\i) at (\i,0.35)
      {\textbf{\si}\\[-1pt]{\scriptsize 10\textsuperscript{\dix} octets}};
    \node[bloc=white, minimum width=17mm, anchor=north] (b\i) at (\i,-0.35)
      {\textbf{\iec}\\[-1pt]{\scriptsize 2\textsuperscript{\deux} octets}};
    \node[et=Pbrun, below=1.5mm of b\i] (p\i) {+\ecart\,\%};
  }
  \node[leg, anchor=east, font=\sffamily\scriptsize\bfseries] at (h1 -| 0.45,0) {décimal (SI)};
  \node[leg, anchor=east, font=\sffamily\scriptsize\bfseries] at (b1 -| 0.45,0) {binaire (CEI)};
  \node[leg, anchor=east, font=\sffamily\scriptsize\bfseries] at (p1 -| 0.45,0) {écart};
\end{tikzpicture}
\end{document}
